% appendices/deployment_manual.tex
% Apéndice B: Manual de instalación y despliegue
% ============================================================================
%
% GUÍA: Manual paso a paso para desplegar el sistema en un entorno local.
%       Los tribunales valoran mucho que el sistema sea reproducible.
%       No se evalúa como contenido del TFG, pero demuestra profesionalidad.
%
%   \section{Requisitos previos}
%   - Docker y Docker Compose instalados (versiones mínimas).
%   - Git para clonar el repositorio.
%   - Sistema operativo: Linux, macOS o Windows con WSL2.
%   - Puertos necesarios libres (8000, 5432, 6379, 9000, etc.).
%
%   \section{Obtención del código fuente}
%   - Clonar el repositorio: git clone ...
%   - Estructura de directorios del proyecto.
%
%   \section{Configuración de variables de entorno}
%   - Copiar .env.example a .env
%   - Describir cada variable de entorno:
%     * DATABASE_URL, REDIS_URL, MINIO_*, SMTP_*, SECRET_KEY, etc.
%   - Ejemplo de .env mínimo funcional.
%
%   \section{Despliegue con Docker Compose}
%   - docker-compose up -d
%   - Verificar que todos los servicios arrancan.
%   - Ejecutar migraciones: docker-compose exec django python manage.py migrate
%   - Crear superusuario: docker-compose exec django python manage.py createsuperuser
%   - Cargar datos iniciales (si hay fixtures).
%
%   \section{Verificación del despliegue}
%   - Acceder a la API: http://localhost:8000/api/
%   - Acceder a Swagger UI: http://localhost:8000/api/docs/
%   - Acceder a MinIO console: http://localhost:9001/
%   - Ejecutar pruebas: docker-compose exec django pytest
%
%   \section{Problemas frecuentes}
%   - Puerto ocupado → cómo cambiar puertos.
%   - Permisos de Docker → sudo o grupo docker.
%   - MinIO no arranca → verificar volúmenes.
% ============================================================================

% TODO: Redactar el manual de despliegue cuando el sistema esté funcional.

[Manual de instalación y despliegue pendiente de redacción.]
